\section{Introduction}\label{sec:introduction}

Ricci-flat metrics can degenerate by developing regions at several widely
separated length scales. In four dimensions, an asymptotically locally
Euclidean end of one region can join a neighborhood of an orbifold point
in another. Repeating this configuration produces a finite rooted tree.
The geometry on every fixed vertex region may converge smoothly while
ratios of scales tend to zero. An estimate useful in this situation must
control the whole joined metric, including the regions between its
limiting pieces.

We study a renormalized Einstein--Hilbert functional on such metrics.
The leading Euclidean divergence in the scalar-curvature density is
subtracted at the exterior end of the root. The resulting quantity,
denoted by $E$, has length degree two. Its differential is the integral
pairing with the Einstein tensor $\tfrac12 Rg-\Ric$.
A fixed root with decay $g-\delta=O(r^{-q})$, $q>1$, makes the
quadratic remainder and the first-variation pairing integrable. The
individual scalar integral and boundary flux need not converge.

Renormalization of curvature energies on noncompact spaces has several
established forms. Haslhofer's renormalized Perelman functional treats
relative energies near steady solitons~\cite{Haslhofer2011}.
Deruelle and Ozuch subtract the mass contribution in their analytic
ALE functional and discuss explicitly the corresponding renormalized
Hilbert--Einstein expression~\cite[Definition~4.2, Remark~4.3, and
Proposition~4.4]{DeruelleOzuch2025}. Their
weighted \L{}ojasiewicz inequality near a fixed Ricci-flat ALE background
holds without an integrability hypothesis; integrability gives an
improved exponent~\cite[Theorem~1.3]{DeruelleOzuch2025}. The functional
$E$ here is the unminimized scalar-curvature-minus-flux expression,
whereas their $\lambda_{\mathrm{ALE}}$ also minimizes over a potential.
Our estimate concerns simultaneous degeneration of the joining scales
in a fixed tree, outside one fixed weighted neighborhood of a smooth
ALE metric.

The main theorem, stated with the geometric hypotheses in
Section~\ref{sec:elliptic}, gives a strict improvement over the direct
first-variation estimate. Let $M$ control ten scaled derivatives of
the dimensionless Ricci tensor, with positive exponential decay on
joining half-ends and decay $r^{-q}$ at the root. For a fixed limiting
tree and sequences with $M\to0$ and all joining lengths tending to
infinity, Theorem~\ref{thm:tree-inequality} proves $E=o(M)$.
The metrics converge to the specified vertex metrics on compact sets,
and their end bounds are uniform through order twenty. Kernels and
cokernels of the linearized equations are retained in the argument.

This is a sequential statement for one fixed tree. Joining lengths may
diverge at unrelated rates, and the proof does not require the
infinitesimal Ricci-flat deformations to integrate to a smooth moduli
space. These features are relevant to four-dimensional Ricci-flow
concentration, where the tree is obtained only after choosing a sequence
of times. For a compact four-dimensional Ricci flow on a finite time interval
with bounded scalar curvature, a curvature blow-up that converges
smoothly to an ancient flow has a Ricci-flat ALE limit
\cite[arXiv version, Corollary~1.9]{BamlerZhangI}.
Related distance-control and compactness results appear
in~\cite{BamlerZhangII,BamlerConvergence}. The present article proves
the static estimate; its hypotheses are stated independently of a flow.

\subsection{Why an analytic path enters the proof}

On a fixed finite-dimensional analytic family, a gradient inequality
relates a function to its differential near a critical set. The
parameters here include neck lengths $L_e\to\infty$. Their boundary
is outside an ordinary analytic coordinate neighborhood, and a change
such as $t_e=e^{-L_e}$ need not make the family analytic at $t_e=0$.
Asymptotic power-logarithm expansions can exist to every finite order
without defining a convergent power series.

Our substitute is a path-selection statement for a precise class of
holomorphic functions. In each length these functions have expansions
in $e^{-dL}$ with polynomial coefficients in $L$, where the
nonnegative exponent sets are locally finite. The coefficients admit
the same expansions in the other lengths. Remainders are strictly
better than the requested exponential order, and all recursively
extracted coefficients are holomorphic on a common ordinary-parameter
neighborhood. These hypotheses are defined in
Definition~\ref{path:def-admissible}.

Theorem~\ref{thm:path-selection} realizes finite sign conditions and
finite prescribed limits along a real analytic path. Every actual
length coordinate and every bounded ordinary coordinate is eventually
monotone or constant. Therefore the path has finite variation in the
weighted differential form
\[
 \sum_\alpha |dz_\alpha|+\sum_e e^{-cL_e}|dL_e|,
 \qquad c>0.
\]
This finite variation is the exact property used in the elliptic
argument. The construction does not replace an all-order expansion
by a convergent series.

Quasianalytic asymptotic algebras provide the ancestry of this approach.
Speissegger~\cite{SpeisseggerQA} develops one-variable Ilyashenko
algebras, and Galal, Kaiser, and Speissegger~\cite{GalalKaiserSpeissegger}
construct algebras based on transserial asymptotic expansions. Their
one-variable field and differentiation results provide precedents for
the clock calculus used here. Our simultaneous length specialization, ordinary
parameter normalization, and path selection are proved in
Section~\ref{sec:path}; no multivariable selection theorem is imported
from those works.

\subsection{The elliptic reduction and its quantitative output}

The analytic family is built in a gauge whose principal part is a
componentwise scalar elliptic operator. The cancellation underlying
this use of coordinates is part of the Ricci-coordinate method;
see DeTurck and Kazdan~\cite{DeTurckKazdan1981} for its classical
geometric setting. On each fixed vertex, finite core tests separate
the kernel, while finitely many compactly supported sources represent
the cokernel. On a joined edge, the two Cauchy matching conditions are
supplemented by the value in the critical spherical channel. Explicit
radial mode matrices then give an inverse whose bound is uniform as
the lengths diverge.

This reduction has close predecessors in Einstein desingularization.
Biquard constructs metrics modulo a finite-dimensional obstruction
space in the Eguchi--Hanson gluing problem
\cite[arXiv version, Proposition~8.1]{Biquard2011}.
Ozuch develops uniform inverse estimates and nearby
Einstein-modulo-obstructions families for desingularization trees,
including ALE orbifolds
\cite[arXiv version, Proposition~11 and Theorem~4.6]{Ozuch2022II}.
The original development was also informed by Ozuch's treatment of
potentially nonintegrable Einstein deformations
\cite{Ozuch2021Integrability}. This is methodological motivation, not
an invocation of that paper's results for the tree inequality proved here.
Our proof retains that obstruction-space viewpoint. The additional
interfaces needed here are continuation in complex length regions,
recursive finite-order expansions with a common ordinary-parameter
domain, and a differential estimate for the renormalized observable.

The ordinary parameters $z$ record these finite data. The length
parameters are continued to quadratic complex regions. A uniform
contraction there gives a small family $h_*(L,z)$ and, for a metric
$g$ in the theorem, the comparison
\[
 \|g-h_*\|\le CM,\qquad
 \|\Ric(h_*)\|\le CM,\qquad
 |E(g)-E(h_*)|\le CM^2.
\]
The norms and the gauge are specified in Section~\ref{sec:family}.
This comparison is why the finite-dimensional family is sufficient
even though the original metrics depend on infinitely many variables.

Set $F=E(h_*)$. Let $J$ be the nonnegative squared Ricci norm defined
in Equation~\eqref{ell:observables}, with an additional weight on the
root tail. The first variation and the shrinking physical volumes of
descendant regions imply
\[
 |dF|\le C\sqrt J
 \left(\sum_\alpha|dz_\alpha|+\sum_e e^{-cL_e}|dL_e|\right).
\]
The remaining analytic obligation is to place $F$ and $J$ in the
path-selection class. We establish this by opening one edge at a
time. Solutions on an opened end have integer exponential tails.
Finite incoming modes cancel the seam mismatch, and triangular linear
auxiliary equations produce each additional asymptotic coefficient.
Finite-part integration transfers the expansions to $F$ and $J$.
Every order uses the same small nonlinear branch and the same basic
inverse; coefficients with greater prescribed growth enter only
linear solves. This preserves one ordinary-parameter domain through
all orders.

If the tree inequality failed, the comparison would produce a sequence
with $F\ne0$, $F\to0$, and $F^2\ge c_0J$ for some $c_0>0$.
Path selection would preserve those conditions on a path of finite
weighted variation. Dividing the differential estimate by $|F|$
would then give finite variation of $\log|F|$, contradicting
$F\to0$. Corollary~\ref{path:variation} isolates this short final
argument from the substantial analytic work needed to apply it.

\subsection{Organization and conventions}

Section~\ref{sec:elliptic} defines the tree geometry, the Ricci error,
and the functional and states the main theorem.
Section~\ref{sec:path} proves path selection.
Section~\ref{sec:family} constructs the inverses, gauge, analytic family,
and comparison with the original sequence.
Sections~\ref{ell:tail-calculus} and~\ref{ell:length-expansion}
prove the end and length expansions.
Section~\ref{ell:scalar-calculus} establishes the scalar expansion
class and completes the theorem.

We use $\Delta=\tr\nabla^2$. Metric and tensor components on a quotient
end are equivariant Cartesian components on its Euclidean cover;
all integrals are quotient integrals. Scaled derivative norms insert
one local length factor per derivative. Tensor magnitudes are physical
unless vertex units are explicitly specified. Constants may depend on
the fixed tree, its charts, and a fixed finite differentiation order.
Uniformity in diverging lengths or in an approximation index is
asserted at the point where it is needed.
